% W4i46_Observables.tex
% DFD 17-Agent Research Programme --- Iteration 46 (i46)
% Agents 8 (Laboratory), 9 (Particle Physics), 10 (Astrophysics), 11 (CMB), 12 (LSS), 13 (Dark Sector), 14 (SymBoltz/Numerics)
% Theme: Phase 0A execution plan, birefringence Chern-Simons assessment, Euclid pre-registration, JUNO first results
% Date: 2026-03-23

\documentclass[11pt,a4paper]{article}
\usepackage[utf8]{inputenc}
\usepackage{amsmath,amssymb,amsfonts,amsthm}
\usepackage{hyperref}
\usepackage{booktabs}
\usepackage{longtable}
\usepackage{geometry}
\usepackage{xcolor}
\usepackage{tcolorbox}
\geometry{margin=2.5cm}

\newtheorem{theorem}{Theorem}
\newtheorem{proposition}{Proposition}
\newtheorem{lemma}{Lemma}
\newtheorem{corollary}{Corollary}
\newtheorem{definition}{Definition}
\newtheorem{remark}{Remark}

\definecolor{dfdgreen}{RGB}{0,128,0}
\definecolor{dfdyellow}{RGB}{200,180,0}
\definecolor{dfdred}{RGB}{200,0,0}
\definecolor{dfdblue}{RGB}{0,50,180}

\newcommand{\GREEN}{\textcolor{dfdgreen}{\textbf{GREEN}}}
\newcommand{\YELLOW}{\textcolor{dfdyellow}{\textbf{YELLOW}}}
\newcommand{\RED}{\textcolor{dfdred}{\textbf{RED}}}
\newcommand{\NEW}{\textcolor{dfdblue}{\textbf{NEW}}}
\newcommand{\RESOLVED}{\textcolor{dfdgreen}{\textbf{RESOLVED}}}
\newcommand{\Tr}{\operatorname{Tr}}
\newcommand{\CP}{\mathbb{C}P}

\title{DFD i46: Observables --- Laboratory, Particle, Astro, CMB, LSS, Dark Sector, Numerics\\[6pt]
\large Agents 8, 9, 10, 11, 12, 13, 14\\[6pt]
\normalsize Iteration 15/20 of Quantum Gravity Cycle (i32--i51)}
\author{DFD 17-Agent Research Programme}
\date{2026-03-23}

\begin{document}
\maketitle
\tableofcontents
\newpage

%=============================================================================
\part{Agent 8: Laboratory Experiments}
\label{part:agent8}
%=============================================================================

\section{Mandate}

Update Th-229 nuclear clock progress. SRF cavities. SQMS. Track $\dot{\alpha}/\alpha$ timeline.

\section{\NEW: Th-229 Frequency Reproducibility --- Nature Published}

\begin{tcolorbox}[colback=green!5!white, colframe=green!60!black,
  title={Nature (January 2026): Solid-State Th-229 Clock Frequency Reproducibility}]

The definitive Th-229 frequency reproducibility result is now published in Nature:
\begin{itemize}
\item At 195~K: reproducibility 220~Hz ($1.1 \times 10^{-13}$) for two differently doped $^{229}$Th:CaF$_2$ crystals over 7 months.
\item Optimal working temperature: $196 \pm 5$~K, where first-order thermal sensitivity vanishes.
\item In-situ temperature co-sensing reduces temperature-induced systematic shift below $10^{-18}$ fractional frequency uncertainty.
\item Frequency: $2020407384335 \pm 2$ kHz (148.382182883 nm).
\end{itemize}

\textbf{DFD significance}: The path to nuclear clock sensitivity of $\dot{\alpha}/\alpha \sim 10^{-19}$/yr by $\sim$2030 is concretely on track. The $10^{-18}$ systematic floor demonstrated at optimal temperature means the nuclear clock can in principle reach DFD's predicted level. The limiting factor is now statistics (integration time), not systematics.

\textbf{Timeline}: At 220~Hz linewidth, reaching $10^{-19}$ fractional precision requires $\sim 10^4$ seconds of integration per measurement. A multi-crystal array (solid-state advantage: orders of magnitude more emitters than atomic clocks) could accelerate this substantially.
\end{tcolorbox}

\section{\NEW: Dark Matter Detection via Nuclear Clock}

\begin{tcolorbox}[colback=blue!5!white, colframe=blue!60!black,
  title={Dark Side of Time: Th-229 Nuclear Clock for Dark Matter Detection (July 2025)}]

A new method proposes using the Th-229 nuclear clock to detect dark matter through variations in the fine-structure constant. The nuclear transition's enhanced sensitivity to $\alpha$ variations ($K_\alpha$ coefficient confirmed) makes it potentially sensitive to oscillating dark matter fields.

\textbf{DFD relevance}: If DFD's $\psi$-field has oscillating components at laboratory frequencies (not predicted by DFD, but not excluded), the nuclear clock could detect them. More importantly, the method validates the $K_\alpha$ sensitivity calculation that underpins DFD's $\dot{\alpha}/\alpha$ prediction.
\end{tcolorbox}

\section{SRF and SQMS: No Change}

\begin{itemize}
\item Dark SRF refined bounds (PRL, March 2026) stand as the latest result.
\item SQMS 2.0 (\$125M renewal) continues.
\item Nb$_3$Sn cavity quality improvement at elevated temperature (2026) --- ongoing.
\end{itemize}

\section{Agent 8 Assessment: Grade A-}

Unchanged from i45. The Th-229 programme remains on track with Nature publication of reproducibility confirmed. The dark matter detection proposal adds a new science case for the nuclear clock.

\section{New Literature (Agent 8)}

\begin{enumerate}
\item \textbf{Nature (January 2026)} --- Th-229 solid-state clock frequency reproducibility. \NEW{} (status: published).
\item \textbf{PhysOrg (July 2025)} --- Dark matter detection via Th-229 nuclear clock method. \NEW.
\item \textbf{Nature Communications (2025)} --- $K_\alpha$ sensitivity coefficient. (continued)
\item \textbf{Current Progress on $^{229}$Th Nuclear Clock} (Photonics 13, 2026) --- Review. (continued)
\item \textbf{Dark SRF refined bounds} (PRL, March 2026) --- (continued)
\item \textbf{SQMS 2.0} (\$125M, 2025--2026) --- (continued)
\end{enumerate}


%=============================================================================
\part{Agent 9: Particle Physics}
\label{part:agent9}
%=============================================================================

\section{Mandate}

LHC Run 3 final months. Moriond 2026 follow-up. HL-LHC preparations.

\section{Moriond 2026: Complete Summary}

\begin{tcolorbox}[colback=blue!5!white, colframe=blue!60!black,
  title={ATLAS Moriond 2026: 45 Analyses --- Full Assessment}]

The 60th anniversary Rencontres de Moriond (15--22 March 2026) results are now fully assessed. Key highlights:

\begin{enumerate}
\item \textbf{$H \to \mu\mu$}: Evidence at $3.4\sigma$ observed ($2.5\sigma$ expected). First evidence in ATLAS.
\item \textbf{$H \to Z\gamma$}: Observed excess of $2.5\sigma$ ($1.9\sigma$ expected) combining Run 2 + Run 3 data.
\item \textbf{$H \to b\bar{b}$ (VBF)}: Signal strength consistent with SM.
\item \textbf{$H \to c\bar{c}$}: Search presented, no excess.
\item \textbf{$HH$ production}: Upper limits from $t\bar{t}HH$ channel.
\item \textbf{SUSY}: Record-setting exclusion limits for higgsinos, long-lived charginos, $\tilde{\tau}$-sleptons, displaced decays, and direct top squark pair production. No signals.
\item \textbf{Top quark}: High-precision measurements.
\end{enumerate}

\textbf{DFD assessment}: All results consistent with SM. No BSM signals. The continued absence of SUSY and other BSM signals strengthens DFD's prediction that the SM particle content is complete. The $H \to \mu\mu$ evidence at $3.4\sigma$ is consistent with DFD's predicted Yukawa structure.

\textbf{$H \to Z\gamma$ note}: The $2.5\sigma$ excess is consistent with the SM one-loop prediction. DFD predicts no enhancement. If this rises to $> 5\sigma$ with more data, it confirms the SM loop structure, which DFD inherits.
\end{tcolorbox}

\section{LHC Run 3: Final 3 Months}

\begin{itemize}
\item Run 3 extends to end of June 2026 (confirmed --- 3 months remaining).
\item LS3 begins July 2026 (confirmed).
\item Integrated luminosity: $\sim$250 fb$^{-1}$ expected by end of Run 3.
\item HL-LHC full-scale magnet tests: $\sim$2/3 of MQXFB quadrupole coils completed.
\item HL-LHC Run 4 start: June 2030.
\end{itemize}

\section{Higgs Mass: No Change}

ATLAS: $m_H = 125.11 \pm 0.11$ GeV. CMS: $m_H = 125.35 \pm 0.15$ GeV. DFD compatible. \GREEN.

\section{Agent 9 Assessment: Grade B+}

Unchanged from i45. Moriond 2026 results fully digested with no surprises for DFD.

\section{New Literature (Agent 9)}

\begin{enumerate}
\item \textbf{ATLAS Moriond 2026 summary} (atlas.cern) --- Full assessment including $H \to Z\gamma$. \NEW{} (expanded from i45).
\item \textbf{HiLumi LHC full-scale tests} (Fermilab, February 2026) --- (continued)
\item \textbf{arXiv:2509.22455} --- Higgs boson prospects at the LHC. (continued)
\item \textbf{ATLAS Higgs mass} (2025) --- $m_H = 125.11 \pm 0.11$ GeV. (continued)
\item \textbf{CMS Higgs mass} (2025) --- $m_H = 125.35 \pm 0.15$ GeV. (continued)
\item \textbf{HL-LHC schedule} (CERN, 2026) --- LS3 from July 2026. (continued)
\end{enumerate}


%=============================================================================
\part{Agent 10: Astrophysics}
\label{part:agent10}
%=============================================================================

\section{Mandate}

Wide binary update. MOND tests. Gaia DR4 timeline.

\section{Wide Binaries: Quality Framework Holds}

No new wide binary results since i45. The quality framework paper (MNRAS 547, 2026) and its application (arXiv:2602.24035) remain the methodological standard. The field continues to converge toward GR with improving methodology.

\section{\NEW: Gaia DR4 Timeline Clarified --- December 2026}

\begin{tcolorbox}[colback=blue!5!white, colframe=blue!60!black,
  title={Gaia DR4: No Earlier Than December 2026}]

The ESA Gaia data release scenario now confirms:
\begin{itemize}
\item \textbf{Gaia DR4}: No earlier than December 2026.
\item DR4 will include epoch-level astrometry for all sources.
\item This enables self-consistent fitting of astrometric data with Keplerian orbital solutions.
\item DR4 will be transformative for binary star research.
\end{itemize}

\textbf{DFD significance}: DR4 with orbital solutions is the definitive wide binary test. DFD predicts no MOND-like departure at any separation. The December 2026 date means the wide binary EFE prediction (\YELLOW) could be resolved within the i32--i51 cycle.

\textbf{Assessment}: EFE prediction remains \YELLOW{} pending DR4. Timeline now concrete: resolution possible by early 2027.
\end{tcolorbox}

\section{Agent 10 Assessment: Grade B+}

Unchanged from i45. DR4 timeline clarification is useful but no new data has arrived.

\section{New Literature (Agent 10)}

\begin{enumerate}
\item \textbf{ESA Gaia DR4 scenario} (cosmos.esa.int) --- December 2026 target. \NEW{} (timeline update).
\item \textbf{arXiv:2510.17746} (October 2025) --- Machine learning identification of Gaia astrometric exoplanet orbits. \NEW.
\item \textbf{arXiv:2602.24035} (February 2026) --- Quality framework applied. (continued)
\item \textbf{Quality framework} (MNRAS 547, 2026) --- Original paper. (continued)
\item \textbf{Cookson et al.\ (2026)}, arXiv:2504.07569 --- GR preferred. (continued)
\item \textbf{Hernandez et al.\ (2024)} --- Critical review. (continued)
\end{enumerate}


%=============================================================================
\part{Agent 11: CMB}
\label{part:agent11}
%=============================================================================

\section{Mandate}

$A_L$ status. Simons Observatory update. \textbf{i46 CRITICAL}: Birefringence risk assessment --- can DFD accommodate $\beta \neq 0$ via Chern--Simons extension?

\section{$A_L$: Unchanged}

$A_L^{\text{recon}} = 1.025 \pm 0.017$ (Qu et al., PRL 136, 2026). DFD prediction $A_L = 1.00$--$1.05$ consistent. Status: \YELLOW{} pending Phase 0.

\section{\NEW: Birefringence Risk Assessment --- Chern--Simons Extension Analysis}

\begin{tcolorbox}[colback=red!5!white, colframe=red!60!black,
  title={CRITICAL: Can DFD Accommodate $\beta \neq 0$?}]

\textbf{The question}: If cosmic birefringence is confirmed at $\beta \sim 0.3^\circ$ or larger, can DFD accommodate it without abandoning the scalar nature of $\psi$?

\textbf{Analysis}:

\underline{Option 1: $\psi$ is purely scalar --- birefringence from independent source.}
\begin{itemize}
\item DFD's $\psi$-field is a scalar. Scalars do not rotate polarization.
\item Birefringence could arise from independent physics: primordial magnetic fields, axion-like particles (ALPs), or Chern--Simons modifications of gravity.
\item This is the conservative option: DFD remains unchanged, birefringence is attributed to a separate sector.
\item \textbf{Problem}: This requires two beyond-SM fields ($\psi$ + ALP/pseudoscalar), reducing DFD's explanatory economy.
\end{itemize}

\underline{Option 2: Chern--Simons coupling of $\psi$ to $F \tilde{F}$.}
\begin{itemize}
\item A scalar $\psi$ can couple to electromagnetism via the Chern--Simons-like term $\mathcal{L} \supset g_{\psi\gamma} \psi F_{\mu\nu} \tilde{F}^{\mu\nu}$.
\item This coupling produces cosmic birefringence if $\psi$ evolves cosmologically.
\item \textbf{However}: This coupling is parity-violating, and $\psi$ is a true scalar (not a pseudoscalar). The coupling $\psi F \tilde{F}$ violates CPT if $\psi$ is a true scalar.
\item DFD's $\psi$ is defined as a scalar refractive field. Making it couple to $F \tilde{F}$ would require either (a) $\psi$ has a pseudoscalar component, or (b) the coupling arises through a composite operator.
\end{itemize}

\underline{Option 3: DFD Chern--Simons sector already exists.}
\begin{itemize}
\item DFD's $\alpha$ derivation uses the Chern--Simons level sum on $S^3$. The CS levels are topological invariants of the gauge bundle.
\item If the CS sector that determines $\alpha$ also generates a $\theta$-parameter for the photon, then a cosmological evolution of the CS sector could produce birefringence.
\item This is the most elegant option: birefringence emerges from the same CS structure that gives $\alpha$.
\item \textbf{Requirement}: The CS level sum must have a time-dependent component that rotates the photon polarization by $\beta \sim \Delta\theta/(4\pi)$.
\end{itemize}

\textbf{Quantitative estimate for Option 3}:

For $\beta = 0.3^\circ = 5.2 \times 10^{-3}$ rad, we need $\Delta\theta \sim 4\pi \beta \sim 0.065$. The Chern--Simons level sum gives $\alpha^{-1} = 137.036$, which involves $\theta = 2\pi/k$ for each level $k$. A cosmological shift $\delta k / k \sim 10^{-4}$ would produce the observed birefringence.

\textbf{But}: The CS levels are integers. They cannot shift continuously. A continuous birefringence signal would require a non-topological (dynamical) component of the CS coupling, which goes beyond DFD's current formalism.

\textbf{Assessment}:
\begin{center}
\begin{tabular}{lcc}
\toprule
\textbf{Option} & \textbf{Compatibility} & \textbf{Economy} \\
\midrule
1. Independent source & Full & Poor \\
2. $\psi F \tilde{F}$ coupling & Requires $\psi$ extension & Medium \\
3. CS sector evolution & Elegant but requires non-integer CS & Best \\
\bottomrule
\end{tabular}
\end{center}

\textbf{Conclusion}: DFD can accommodate birefringence through Option 1 (trivially) or Option 3 (elegantly, if the CS formalism is extended to include non-topological dynamical terms). Option 2 is the least compatible because it requires changing $\psi$'s scalar nature.

\textbf{Risk level}: Remains \textbf{MEDIUM-HIGH}. The birefringence signal is not yet confirmed at $5\sigma$. Current data: $\beta = (0.30 \pm 0.11)^\circ$ from Planck PR4, with new uncertainty treatment suggesting it may be larger.

\textbf{Next steps}: Track anisotropic birefringence constraints (arXiv:2504.13154) and simulation-based inference (arXiv:2602.12019). SO SAT data (deploying 2026) will be decisive.
\end{tcolorbox}

\section{\NEW: Anisotropic Birefringence Constraints}

\begin{tcolorbox}[colback=blue!5!white, colframe=blue!60!black,
  title={arXiv:2504.13154 (June 2025): Anisotropic Birefringence from CMB B-mode}]

Joint analysis using SPTpol + ACT + POLARBEAR + BICEP polarization data yields best-fit anisotropic birefringence amplitude $0.42^{+0.40}_{-0.34} \times 10^{-4}$, \textbf{consistent with zero}.

\textbf{DFD assessment}: The anisotropic component is not detected. This is important: if birefringence is isotropic only, it is more consistent with a homogeneous scalar field evolution (Options 2 or 3) than with localized sources. A purely isotropic signal would be the most natural outcome of any DFD Chern--Simons extension.
\end{tcolorbox}

\section{\NEW: Simulation-Based Inference for Birefringence Origin}

arXiv:2602.12019 (February 2026): Uses simulation-based inference to distinguish dark energy vs.\ dark matter origins for cosmic birefringence. This methodology could eventually distinguish DFD's CS-based birefringence (if it exists) from ALP models.

\section{Simons Observatory: SO:UK SATs Deploying}

\begin{itemize}
\item LAT first light February 22, 2025 (Mars observation).
\item Three SATs operational since mid-2024.
\item SO:UK SATs (90 + 150 GHz dichroic) deploying 2026.
\item NSF \$52M upgrade funds enhanced LAT + 5-year survey through 2033.
\item \NEW{} Polarization angle calibration methodology developed using sparse wire grid (arXiv:2512.19102).
\item \NEW{} PROTOCALC drone-borne calibration source tested in lab.
\end{itemize}

\section{Agent 11 Assessment: Grade B+}

Unchanged from i45. The birefringence Chern--Simons analysis is the major i46 contribution. The risk level remains MEDIUM-HIGH. The anisotropic null result is mildly encouraging for DFD.

\section{New Literature (Agent 11)}

\begin{enumerate}
\item \textbf{arXiv:2504.13154} (June 2025) --- Anisotropic birefringence constraints from B-mode. \NEW.
\item \textbf{arXiv:2602.12019} (February 2026) --- Simulation-based inference for birefringence origin. \NEW.
\item \textbf{arXiv:2512.19102} (December 2025) --- SO SAT polarization angle calibration. \NEW.
\item \textbf{PRL (January 2026)} --- Birefringence angle may be larger. (continued)
\item \textbf{NSF \$52M SO upgrade} (2026). (continued)
\item \textbf{arXiv:2603.12335} (March 2026) --- Higher-order CMB lensing statistics. (continued)
\item \textbf{Qu et al.\ (PRL 136, 2026)} --- $A_L = 1.025 \pm 0.017$. (continued)
\end{enumerate}


%=============================================================================
\part{Agent 12: Large-Scale Structure}
\label{part:agent12}
%=============================================================================

\section{Mandate}

\textbf{i46 CRITICAL}: Euclid pre-registration finalization. Phase 0A-level predictions.

\section{Euclid Pre-Registration: 88 Days --- Concrete Plan}

\begin{tcolorbox}[colback=red!5!white, colframe=red!60!black,
  title={DEADLINE: Euclid Pre-Registration --- 20 June 2026 (88 days)}]

\textbf{Key dates}:
\begin{itemize}
\item \textbf{Euclid Q2}: 24 June 2026 --- Galactic Bulge Survey (NOT cosmology).
\item \textbf{Euclid DR1}: 21 October 2026 --- $\sim$2100 deg$^2$, weak lensing shear, cosmology.
\item \textbf{DFD pre-registration}: Must be submitted before Q2 (20 June) to establish timestamp.
\end{itemize}

\textbf{Phase 0A plan for pre-registration content}:

\begin{enumerate}
\item \textbf{Distance-bias prediction}: $D_L^{\text{DFD}}(z) = D_L^{\Lambda\text{CDM}}(z) \cdot e^{\Delta\psi(z)}$, with $\Delta\psi(z) = \sum_{n=1}^{3} a_n z^n$.
\item \textbf{Qualitative direction}: DFD predicts a modification to the weak lensing convergence power spectrum $C_\ell^{\kappa\kappa}$ that shifts the inferred $(\Omega_m, \sigma_8)$ toward lower $S_8$ relative to CMB-only.
\item \textbf{$A_L$ background-only prediction}: From Phase 0A, estimate $A_L^{\text{DFD}} \in [1.00, 1.05]$ (preliminary numerical range, consistent with current data).
\item \textbf{BAO distance ratios}: DFD-modified $D_V(z)/r_d$ at Euclid redshift bins.
\item \textbf{Null hypothesis}: $\Lambda$CDM. Alternative: DFD distance-bias.
\item \textbf{Statistical threshold}: $3\sigma$ in blind analysis of DR1 shear catalogue.
\end{enumerate}

\textbf{Contingency}: If Phase 0A cannot produce numerical BAO predictions by 20 June, submit qualitative pre-registration with the distance-bias formula, predicted direction of deviation, and explicit statement that numerical predictions will follow in a separate paper.
\end{tcolorbox}

\section{\NEW: Euclid Higher-Order WL Statistics --- Fisher Forecasts}

\begin{tcolorbox}[colback=blue!5!white, colframe=blue!60!black,
  title={Euclid Prep.\ LXXXV (A\&A, March 2026): DR1 Higher-Order WL Statistics}]

The Euclid preparation paper (arXiv:2510.04953, published A\&A March 2026) provides Fisher forecasts for five higher-order weak lensing statistics (PDF, $\ell_1$-norm, peak counts, Minkowski functionals, Betti numbers) at DR1 depth:

\begin{itemize}
\item Each higher-order statistic outperforms shear two-point correlation functions by a factor of $\sim 2.5$ on $w_0$ constraints.
\item LensMC is the designated shape measurement method for DR1.
\item Tomographic analysis with realistic photometric redshifts.
\end{itemize}

\textbf{DFD relevance}: Higher-order statistics are more sensitive to DFD's distance-bias than two-point correlations. The factor $2.5\times$ improvement means Euclid DR1 could detect DFD effects at $\sim 2.5$ times the significance of the shear power spectrum alone. This strengthens the case for pre-registration.

\textbf{For Phase 0B}: If the $\psi$ perturbation theory is worked out, DFD predictions for non-Gaussian convergence statistics (peak counts, Minkowski functionals) would be a powerful discriminant.
\end{tcolorbox}

\section{\NEW: Euclid Early Release Observations --- WL Pipeline Validated}

Euclid Early Release Observations (arXiv:2507.07629) demonstrated weak lensing measurement capabilities on real data, including combined strong and weak lensing solutions for Abell 2390 extending beyond its virial radius. This validates the WL pipeline that will be applied to DR1.

\section{DESI DR2: Dynamical DE Debate Continues}

No change from i45. Giar\`e et al.\ (EPJC, 2025) challenge stands. DFD's distance-bias provides a geometric alternative to $w_0 w_a$.

\section{Agent 12 Assessment: Grade B+}

Unchanged from i45. The pre-registration plan is now more concrete with the higher-order WL statistics Fisher forecasts. The 88-day countdown is active.

\section{New Literature (Agent 12)}

\begin{enumerate}
\item \textbf{Euclid Prep.\ LXXXV} (A\&A, March 2026; arXiv:2510.04953) --- DR1 higher-order WL Fisher forecasts. \NEW{} (published).
\item \textbf{arXiv:2507.07629} --- Euclid ERO weak lensing pipeline validation. \NEW.
\item \textbf{ESA Euclid DR1} (cosmos.esa.int) --- DR1 content and timeline. \NEW{} (status update).
\item \textbf{Giar\`e et al.\ (EPJC, 2025)}, arXiv:2504.15222 --- Challenges DESI DR2 DE. (continued)
\item \textbf{DESI DR2 BAO} (Phys.\ Rev.\ D) --- (continued)
\item \textbf{SymBoltz.jl} (A\&A, March 2026) --- (continued)
\end{enumerate}


%=============================================================================
\part{Agent 13: Dark Sector}
\label{part:agent13}
%=============================================================================

\section{Mandate}

$\Sigma m_\nu$ monitoring. JUNO timeline and first results. Dark sector constraints.

\section{\NEW: JUNO First Results --- World-Leading Precision}

\begin{tcolorbox}[colback=green!5!white, colframe=green!60!black,
  title={JUNO First Physics Results (November 2025): $\Delta m^2_{21}$ and $\theta_{12}$}]

Major milestone: Using 59.1 days of data (26 August -- 2 November 2025), JUNO achieved:
\begin{itemize}
\item World-leading precision on $\Delta m^2_{21}$ --- 1.6$\times$ improvement over all previous measurements combined.
\item Simultaneous high-precision determination of $\theta_{12}$.
\item Percent-level precision in atmospheric mass splitting expected within months of additional data.
\end{itemize}

\textbf{Mass ordering status}: Combination of JUNO + T2K + NOvA gives slight preference for \textbf{normal ordering} at $2.2$--$2.3\sigma$ (p-value for IO: 2--2.6\%).

\textbf{DFD assessment}: DFD predicts normal ordering. The $2.2\sigma$ preference is encouraging but not decisive. The $3\sigma$ milestone is expected by 2027--2028 with $\sim$1 year of JUNO data combined with T2K/NOvA.
\end{tcolorbox}

\section{$\Sigma m_\nu$: Marginal but Alive (Updated)}

\begin{itemize}
\item $\Sigma m_\nu < 0.064$ eV (DESI DR2 + Planck, $\Lambda$CDM).
\item $\Sigma m_\nu < 0.053$ eV (Feldman--Cousins approach) --- \textbf{below oscillation floor}.
\item $\Sigma m_\nu < 0.082$ eV (ACT DR6 + DESI DR2, $w_0 w_a$CDM).
\item DFD: 0.061 eV.
\end{itemize}

\textbf{Updated assessment}: The Feldman--Cousins tension ($3\sigma$ between cosmology and oscillations) is a known issue. If resolved by invoking dynamical DE, DFD is safe. If it persists in $\Lambda$CDM, DFD's 0.061 eV remains within the standard 95\% CL bound but is under increasing pressure.

\section{\NEW: Second-Order CPL Models and $\Sigma m_\nu$}

arXiv:2508.16238: Cosmological constraints on neutrino masses in a second-order CPL dark energy model show that additional DE freedom relaxes $\Sigma m_\nu$ bounds to $< 0.10$--$0.12$ eV. DFD's 0.061 eV is comfortably within these extended bounds.

\section{Agent 13 Assessment: Grade B}

Upgrade consideration from B to B+ based on JUNO first results. However, JUNO has not yet reached mass ordering sensitivity, so the upgrade is deferred.

\textbf{Honest note}: Agent 13 remains at B because no DFD-specific computation has been performed. The JUNO results are important context, but Agent 13 should be producing DFD-specific predictions for JUNO sensitivity, not just tracking external results.

\section{New Literature (Agent 13)}

\begin{enumerate}
\item \textbf{arXiv:2511.14593} (November 2025) --- JUNO first reactor neutrino oscillation measurement. \NEW.
\item \textbf{arXiv:2601.09791} (January 2026) --- Lessons from the first JUNO results. \NEW.
\item \textbf{arXiv:2603.17181} (March 2026) --- JUNO + long-baseline synergy. \NEW.
\item \textbf{arXiv:2508.16238} --- Neutrino masses in second-order CPL DE. \NEW.
\item \textbf{arXiv:2503.14744} (March 2025) --- DESI DR2 + Planck $\Sigma m_\nu$. (continued)
\item \textbf{JUNO status} (2026) --- Data taking commenced. (continued)
\end{enumerate}


%=============================================================================
\part{Agent 14: SymBoltz / Numerics --- PHASE 0A EXECUTION PLAN}
\label{part:agent14}
%=============================================================================

\section{Mandate}

\textbf{CRITICAL}: Phase 0A execution plan. Phase 0B $\psi$ perturbation theory derivation begun.

\section{Phase 0A: Concrete Execution Plan}

\begin{tcolorbox}[colback=green!5!white, colframe=green!60!black,
  title={Phase 0A: 12-Hour Background-Only Computation --- Detailed Plan}]

\textbf{Step 1 (Hours 0--2): SymBoltz Installation and Baseline}
\begin{enumerate}
\item Install SymBoltz.jl v1.0.0 (Julia package).
\item Run $\Lambda$CDM baseline with Planck 2018 best-fit parameters.
\item Verify: $C_\ell^{TT}$, $C_\ell^{EE}$, $C_\ell^{\phi\phi}$ match Planck to $< 0.1\%$.
\item Verify: $D_V(z)/r_d$ at DESI fiducial redshifts.
\end{enumerate}

\textbf{Step 2 (Hours 2--5): Distance-Bias Implementation}
\begin{enumerate}
\item Implement custom distance function: $D_L^{\text{DFD}}(z) = D_L^{\Lambda\text{CDM}}(z) \cdot e^{\Delta\psi(z)}$.
\item Parameterize: $\Delta\psi(z) = a_1 z + a_2 z^2 + a_3 z^3$.
\item Fix $a_n$ from DFD field equations at background level:
\begin{itemize}
\item $a_1 = -\frac{1}{2} H_0 \lambda_\psi$ (first-order distance modification).
\item $a_2$ from second-order expansion of FRW $\psi$-field solution.
\item $a_3$ from third-order correction (expected to be small).
\end{itemize}
\item $\lambda_\psi$ is the DFD parameter controlling the $\psi$-field amplitude; bounded by $|\lambda_\psi - 1| < 1.56 \times 10^{-9}$ (SQMS).
\end{enumerate}

\textbf{Step 3 (Hours 5--8): $A_L$ and BAO Computation}
\begin{enumerate}
\item Compute $A_L^{\text{DFD}}$ by evaluating the CMB lensing potential with modified distances.
\item Note: In background-only, $A_L$ changes only through the geometric effect on photon paths, not through perturbation-level changes in the lensing potential.
\item Compute DFD-modified $D_V(z)/r_d$ at Euclid redshift bins.
\item Compare with DESI DR2 measurements.
\end{enumerate}

\textbf{Step 4 (Hours 8--12): Euclid Pre-Registration Numbers}
\begin{enumerate}
\item Generate background-only predictions for Euclid DR1.
\item Compute predicted shift in $(\Omega_m, \sigma_8)$ plane from distance-bias.
\item Estimate Fisher forecast for DFD detectability with DR1 shear catalogue.
\item Write qualitative pre-registration document.
\end{enumerate}

\textbf{Output}: Approximate $A_L$, BAO predictions, qualitative Euclid pre-registration.

\textbf{Limitation}: No perturbation-level predictions. Cannot compute $C_\ell^{TT}$ shape changes, $f\sigma_8(z)$, or ISW effect.
\end{tcolorbox}

\section{Phase 0B: $\psi$ Perturbation Theory --- Derivation Framework}

\begin{tcolorbox}[colback=yellow!5!white, colframe=yellow!60!black,
  title={Phase 0B: $\psi$ Perturbation Theory --- Theoretical Framework (NEW)}]

\textbf{The central question}: What is the perturbation equation for $\delta\psi$ in the cosmological context?

\textbf{Starting point}: DFD defines $\psi$ as a scalar refractive field sourced by matter:
\begin{equation}
\nabla^2 \psi = -4\pi G \rho / c^2 \cdot f(\psi)
\end{equation}
where $f(\psi)$ encodes the nonlinear self-coupling.

\textbf{Perturbed equation} (Newtonian gauge, $\delta\psi = \psi - \bar{\psi}$):
\begin{equation}
\ddot{\delta\psi} + 3H\dot{\delta\psi} + \left(\frac{k^2}{a^2} + m_\psi^2\right)\delta\psi = -\frac{4\pi G}{c^2} \bar{\rho} \delta \cdot f'(\bar{\psi}) + \text{metric perturbation terms}
\end{equation}

\textbf{Key parameters to determine}:
\begin{enumerate}
\item \textbf{$\psi$ mass} $m_\psi$: From the second derivative of the $\psi$ potential. If $m_\psi = 0$, $\psi$ is massless and perturbations propagate at $c$. If $m_\psi \neq 0$, there is a Compton wavelength below which $\psi$ perturbations are suppressed.
\item \textbf{Sound speed} $c_s^2$: For a canonical scalar, $c_s^2 = 1$. If the kinetic term is non-canonical (K-essence-like), $c_s^2$ differs.
\item \textbf{Anisotropic stress} $\pi_\psi$: For a canonical scalar, $\pi_\psi = 0$. This simplifies the Boltzmann hierarchy considerably.
\item \textbf{Coupling to photons and baryons}: Does $\psi$ couple directly to the photon perturbation $\delta_\gamma$? In DFD, $\psi$ modifies the photon propagation via the refractive index, which implies a direct coupling at the perturbation level.
\item \textbf{Initial conditions}: $\delta\psi$ at recombination. If $\psi$ tracks matter, $\delta\psi / \bar{\psi} \approx \delta_m$ initially.
\end{enumerate}

\textbf{Relevant recent literature}:
\begin{itemize}
\item \textbf{arXiv:2602.19576} (February 2026): K-essence Boltzmann dynamics with modified mass-shell condition and geodesic equations. Directly relevant to DFD's modified photon propagation.
\item \textbf{arXiv:2512.16404} (December 2025): Implementing $F(T)$ gravity in Boltzmann codes. Framework for adding modified gravity to SymBoltz.
\item \textbf{SymBoltz.jl} (A\&A, March 2026): Built-in support for quintessence and Brans--Dicke as ``replaceable physical components.'' DFD $\psi$ could be implemented as a custom component.
\end{itemize}

\textbf{SymBoltz.jl capability}: SymBoltz's symbolic-numeric interface allows new species to be added as self-contained submodels. A ``DFD $\psi$'' submodel would specify:
\begin{enumerate}
\item Background equations for $\bar{\psi}(a)$.
\item Perturbation equations for $\delta\psi(k, a)$.
\item Coupling equations to photon and baryon Boltzmann hierarchies.
\end{enumerate}

\textbf{Timeline estimate}: The theoretical derivation (specifying all 5 parameters above) is $\sim$2--4 weeks of focused work. The SymBoltz implementation is $\sim$1--2 weeks after the theory is in hand. Total: 1--2 months for Phase 0B.

\textbf{Status}: Framework specified. Derivation not yet performed.
\end{tcolorbox}

\section{Agent 14 Assessment: Grade B}

Unchanged from i45. Phase 0A plan is now concrete and actionable. Phase 0B theoretical framework is specified. Neither has been executed.

\textbf{Honest note}: Agent 14 has been at the planning stage for five iterations (i42--i46). The plan has improved each iteration (generic $\to$ honest blocker $\to$ 0A/0B split $\to$ concrete steps $\to$ theoretical framework). But planning without execution is not progress. Agent 14 must \emph{execute} Phase 0A in the next available computational session. Grade will remain B until code is run.

\section{New Literature (Agent 14)}

\begin{enumerate}
\item \textbf{arXiv:2602.19576} (February 2026) --- K-essence Boltzmann dynamics. \NEW.
\item \textbf{arXiv:2512.16404} (December 2025) --- $F(T)$ gravity in Boltzmann codes. \NEW.
\item \textbf{SymBoltz.jl} (A\&A, March 2026) --- Published; custom submodel support. (continued, now published)
\item \textbf{arXiv:2509.24740} --- SymBoltz.jl preprint. (continued)
\item \textbf{arXiv:2307.05600} --- Scalar field DM in CLASS. (continued)
\item \textbf{gCAMB} (arXiv:2509.25110) --- (continued)
\end{enumerate}


\end{document}
